%! The Language of Variation — Unit 1, Lesson 1.1 — Experience & Formalize
%! (Activity · QuickNotes · Application · Check Your Understanding)
%! Directory stays `experience/` (a build identifier in shared/lesson.mk); the label
%! students and teachers read is "Experience & Formalize".
% EFFL component, Math Medic sizing: 12pt body, OPEN white space after every question.
% \answerspace{H}{} reserves exactly H in the blank; the key passes the answer as the second
% argument, so the two files paginate identically by construction. Keep the macro
% byte-identical in the blank and the key. \boxguard counts here are 12pt-relative (~14-16).
% Page budget: Activity <=2pp · QuickNotes 1/2pp · Application 1/2-1pp · CYU 1-2pp.
%
% SPOILER RULE: the Activity never says categorical, quantitative, numerical, or measurement
% variable. Students invent their own two piles and name them in their own words ("amounts",
% "which kind"); the debrief attaches the formal terms in QuickNotes. `Individual', `variable',
% and `value' WERE formalized in Lesson 1.0, so they are prior knowledge here, not spoilers.
%
% The \begin{work} block in scenario 2b is authored BYTE-IDENTICALLY in this file and in
% experience_key/. The blank reserves its exact height and prints nothing; the key prints it
% in keyred. That block plus \answerspace is what keeps the two files page-for-page.
\documentclass[12pt]{article}
\usepackage{apstats-article}
\usepackage{apstats-boxes}
\newcommand{\answerspace}[2]{\par\nopagebreak\noindent\begin{minipage}[t][#1][t]{\linewidth}%
  \color{keyred}\bfseries #2\end{minipage}\par}

\begin{document}

\pageheader{Unit 1, Lesson 1.1}{Experience \& Formalize: The Trailhead Register}
% NAMESTRIP: no \namedateperiod here — the name row lives on the cover only.

\vspace{0.02in}

% ── Part 1: the Activity ─────────────────────────────────────────────────────
% Scenario 1 makes the groups INVENT the two piles and name them themselves — no bins are
% handed to them. Scenario 2 breaks the naive rule they will have written ("has digits in it")
% with three digit-written labels and one genuine count, and carries the CRUX in 2b: the
% average of two permit numbers computes fine and means nothing.
\begin{headlinebox}{sky}
\textbf{The Trailhead Register.} Every party that hikes at Ross Hollow State Park signs the
register at the trailhead and signs out again on the way home. Work through the two problems
below with your group using only what you already know, and be ready to explain \emph{how you
know} each answer \textbf{from the register itself}.
\end{headlinebox}

\vspace{0.04in}

\begin{scenariobox}[1. Two Piles]{navy}
\small
\renewcommand{\arraystretch}{1.12}
\begin{tabularx}{\linewidth}{@{} l Y Y Y Y Y @{}}
\rowcolor{sky}
\textbf{Party} & \textbf{Trail} & \textbf{Group size} & \textbf{Time out (min)} &
\textbf{Dog along} & \textbf{Weather} \\
\hline
Alvarez  & Ridge Loop   & 4 &  95 & Yes & Sunny  \\
Boone    & Cedar Falls  & 2 & 140 & No  & Cloudy \\
Chen     & Ridge Loop   & 6 &  80 & Yes & Sunny  \\
Delgado  & Summit Trail & 3 & 215 & No  & Rain   \\
Egan     & Cedar Falls  & 2 & 155 & Yes & Cloudy \\
\hline
\end{tabularx}

\normalsize
\begin{enumerate}[label=\textbf{\alph*.}, leftmargin=1.8em, itemsep=6pt]
  \item One \textbf{row} of this register describes \blank{4.6cm}.
  \item Sort the five recorded headings into \textbf{two piles of your own}. Write each
        heading under the pile where you would put it.

        \begin{minipage}[t]{0.47\linewidth}
        \textbf{Pile 1}
        \answerspace{1.4cm}{}
        \end{minipage}\hfill
        \begin{minipage}[t]{0.47\linewidth}
        \textbf{Pile 2}
        \answerspace{1.4cm}{}
        \end{minipage}
  \item Now name your piles. Finish this sentence \emph{twice}, once for each pile:
        \emph{``A heading belongs in this pile when \ldots{}''}
        \answerspace{2.0cm}{}
  \item For \textbf{one} of your piles, adding up all five entries in a column and dividing by 5
        gives the ranger something useful. For the other it gives nonsense. Say which is which,
        write down the nonsense you would get on \textbf{Weather}, and say what you would do
        with that column instead to summarize it.
        \answerspace{2.4cm}{}
\end{enumerate}
\end{scenariobox}

\boxguard[16]
\begin{scenariobox}[2. Four Columns Written in Digits]{navy}
Next season the ranger adds four more columns. Here are the first three rows.

\vspace{4pt}
\small
\renewcommand{\arraystretch}{1.12}
\begin{tabularx}{\linewidth}{@{} l Y Y Y Y @{}}
\rowcolor{sky}
\textbf{Party} & \textbf{Permit \#} & \textbf{Parking spot} & \textbf{Home ZIP} &
\textbf{Water bottles} \\
\hline
Alvarez & 2041 & 17 & 80211 & 6 \\
Boone   & 6318 & 04 & 80304 & 2 \\
Delgado & 1150 & 11 & 80211 & 3 \\
\hline
\end{tabularx}

\normalsize
\begin{enumerate}[label=\textbf{\alph*.}, leftmargin=1.8em, itemsep=6pt]
  \item Every one of these four columns is written entirely in digits. Put each heading into
        \textbf{one of your two piles} from problem 1.

        \vspace{2pt}

        \textbf{Permit \#} \blank{2.6cm} \hfill \textbf{Parking spot} \blank{2.6cm}

        \vspace{4pt}
        \textbf{Home ZIP} \blank{2.6cm} \hfill \textbf{Water bottles} \blank{2.6cm}
  \item The Alvarez and Boone parties had permit numbers 2041 and 6318, and carried 6 and 2
        water bottles. Compute both averages.
        \begin{work}
        \text{average permit \#} &= \frac{2041 + 6318}{2} = 4179.5 \\[4pt]
        \text{average water bottles} &= \frac{6 + 2}{2} = 4
        \end{work}
        One of those two answers tells the ranger something real about the two parties and one
        does not. Which is which, and what exactly goes wrong with the other?
        \answerspace{1.8cm}{}
  \item So being written in digits is not what decides the pile. Finish the rule in your own
        words: \emph{``A heading belongs in the pile you can average only when \ldots{}''}
        \answerspace{1.6cm}{}
  \item The season report says \textbf{``62\% of parties hiked with a dog.''} That number came
        out of the \emph{Dog along} column. Does it move that column into the average pile?
        \answerspace{2.0cm}{}
\end{enumerate}
\end{scenariobox}

% ── Part 2: QuickNotes (the debrief fills this) — target: half a page ────────
\boxguard[14]
\begin{tcolorbox}[enhanced, colback=sky, colframe=navy, boxrule=0.8pt, arc=2mm,
    left=3mm, right=3mm, top=1.5mm, bottom=1.5mm, breakable,
    fonttitle=\bfseries\small\color{white}, title={QuickNotes: Two Kinds of Variable},
    attach boxed title to top left={yshift=-2mm, xshift=4mm},
    boxed title style={colback=navy, colframe=navy, arc=1.5mm}]
\small
\begin{minipage}[t]{0.34\linewidth}
\vspace{2pt}
\renewcommand{\arraystretch}{1.15}
\begin{tabular}{@{}l l@{}}
\multicolumn{2}{@{}l}{\footnotesize\textbf{\color{navy}Trail} \quad \textbf{\color{navy}Time (min)}}\\
\hline
Ridge Loop   & 95  \\
Cedar Falls  & 140 \\
Summit Trail & 215 \\
\hline
\end{tabular}\\[3pt]
{\footnotesize Left column: \textbf{category}\\
\textbf{names}. Right column:\\
a \textbf{measured amount}.\\
Same five parties, two\\
different kinds of column.}
\end{minipage}\hfill
\begin{minipage}[t]{0.62\linewidth}
\vspace{-4pt}
\begin{itemize}[leftmargin=1.1em, itemsep=5pt]
  \item A \textbf{variable} is a characteristic that \blank{2.2cm} from one individual to
        another.
  \item A \textbf{categorical variable} takes values that are \blank{2.0cm} names or
        group \blank{1.8cm}.
  \item A \textbf{quantitative variable} takes \blank{2.2cm} values for a measured or
        \blank{1.8cm} quantity.
  \item \textbf{The test is not ``does it have digits.''} Ask: is the value an
        \blank{2.0cm} of something? If adding or averaging the values is
        \blank{2.4cm}, it is quantitative. A ZIP code is all digits and is
        \blank{2.6cm}.
  \item Counting a categorical variable \emph{produces} numbers --- counts and percents.
        The \blank{1.9cm} is numeric; the \blank{1.9cm} is still categorical.
\end{itemize}
\end{minipage}
\end{tcolorbox}

% ── Part 3: the Application (worked together, ~6 min) ────────────────────────
\boxguard[14]
\begin{notesbox}{Application: The Season Pass}
We will do this one together. A ski area records five things about each of its 3{,}400
season-pass holders: \textbf{pass number}, \textbf{age} (years), \textbf{pass tier}
(Gold / Silver / Bronze), \textbf{miles traveled} to the mountain, and \textbf{locker number}.

\begin{enumerate}[label=\textbf{\alph*.}, leftmargin=1.8em, itemsep=6pt]
  \item The individuals are \blank{5.6cm}. Label each variable \textbf{C} or \textbf{Q}.

        \vspace{4pt}
        \renewcommand{\arraystretch}{1.3}
        \begin{tabularx}{\linewidth}{@{} Y Y Y Y Y @{}}
        \textbf{Pass \#} & \textbf{Age} & \textbf{Pass tier} & \textbf{Miles traveled} &
        \textbf{Locker \#} \\
        \blank{1.6cm} & \blank{1.6cm} & \blank{1.6cm} & \blank{1.6cm} & \blank{1.6cm} \\
        \end{tabularx}
  \item \textbf{What if we change the variable?} The resort stops recording age in years and
        records \textbf{age group} instead (Under 18 / 18--34 / 35--54 / 55+). Same people,
        same information --- which kind is it now, and what changed?
        \answerspace{1.8cm}{}
  \item A report states: \textbf{``The average pass number is 5{,}120.''} Say in one sentence
        why that number tells you nothing about the pass holders.
        \answerspace{1.6cm}{}
\end{enumerate}
\end{notesbox}

% ── Part 4: Check Your Understanding — practice, NOT scored ──────────────────
\boxguard[14]
\begin{notesbox}{Check Your Understanding}
Work with your partner, then finish on your own. \emph{These are for practice --- they are not
scored.}
\begin{enumerate}[leftmargin=1.9em, itemsep=7pt]
  \item A coffee shop logs every order it fills. Label each variable \textbf{C} or \textbf{Q}.

        \vspace{4pt}
        \renewcommand{\arraystretch}{1.3}
        \begin{tabularx}{\linewidth}{@{} Y Y Y Y @{}}
        \textbf{Drink size (S/M/L)} & \textbf{Number of shots} & \textbf{Price paid (\$)} &
        \textbf{Loyalty card \#} \\
        \blank{1.9cm} & \blank{1.9cm} & \blank{1.9cm} & \blank{1.9cm} \\
        \end{tabularx}

        The individuals here are \blank{5.2cm}.
  \item \textbf{Same survey, two columns.} A pet survey records \emph{Number of pets owned} and
        \emph{Type of pet owned}. Classify each, then say what a summary of each column would
        look like.
        \answerspace{2.0cm}{}
  \item A hospital records each patient's exact body temperature in $^\circ$C. It then adds a
        second column marking each patient \emph{Normal} or \emph{Fever}. Classify both
        columns, and explain where the second column came from.
        \answerspace{2.0cm}{}
  \item Name one variable that is written entirely in digits and is \textbf{not} quantitative,
        and say what its digits are doing instead of measuring.
        \answerspace{1.6cm}{}
  \item \textbf{Formative check.} A survey of high school students records the five variables
        below. Circle the one that is \textbf{quantitative}, then justify it in one sentence.
        \begin{enumerate}[label=\textbf{(\Alph*)}, leftmargin=2.2em, itemsep=3pt]
          \item The number on the student's team jersey
          \item The number of the bus route the student rides
          \item The number of minutes the student spent on homework last night
          \item The student's home ZIP code
          \item The student's cafeteria PIN
        \end{enumerate}
        \answerspace{1.6cm}{}
\end{enumerate}
\end{notesbox}

\end{document}
